\cotefile{33}
\scanpage{98}
supposer $F$ \uline{constant}, puis (\unsure{le résolvant|le résoudre} par un $A_Y^n$) $F=A_Y$, mais alors le th est trivial.

b) Cas $f$ une immersion. Alors le th. est essentiellement trivial, sans hyp. de finitude. En effet, on est ramené au cas d'une imm fermée. Alors
\[
\begin{aligned}
\mathbb R^{\deleted{\unsure{!}}}f^!(\mathbb R\mathscr Hom(F_{*},G^{\bullet}))&\simeq f^!\deleted{\unsure{\mathbb R}}\underbrace{\mathscr Hom^{\bullet}(F_{\bullet},C^{\bullet}(G^{\bullet}))}_{\substack{\text{N.B. c'est flasque,}\\\text{donc }f^!\text{-acyclique}}}\simeq\\
\mathscr Hom^{\bullet}(f^*F_{\bullet},\underbrace{\deleted{\unsure{C}}f^!C^{\bullet}(G^{\bullet})}_{\text{N.B. c'est injectif}})&\simeq\mathbb R\mathscr Hom(\deleted{\unsure{\mathbb R}}\mathbb L f^*(F_{\bullet}),\mathbb Rf^!(G^{\bullet}))
\end{aligned}
\]

\noindent\rule{.35\linewidth}{.3pt}

Les formules de projection \unsure{et de restriction|et d'induction|et de bidirection|d'induction} comme \uline{formules de} \deleted{\unsure{commutation resp.}} \uline{d'échange}
\[
\left\{\begin{aligned}
\mathbb R f_*D_X&\simeq D_Y\mathbb R_!f\\
\deleted{\unsure{\mathbb L}}\mathbb R^!fD_Y&\simeq D_X\mathbb L f^*
\end{aligned}\right.
\]
i.e. $D$ ``échange'' $\mathbb Rf_*$ et $\mathbb R_!f$, $\deleted{\unsure{\mathbb R}}\mathbb Lf^*$ et $\mathbb L^!f$.
% archive sheet number: 97
\scanpage{99}
\mspage{72}
% archive sheet number: 98
\[
(\overset{\handwritten{\sim}}{F}_x(u,v))(a_1,\ldots,a_p,b_1,\ldots,b_p)=(F_x(u-v))(a_1-u,\ldots,a_p-u,b_1-v,\ldots,b_p-v)
\]
where $a_i,b_i\handwritten{\in}T_x$. We denote by $\handwritten{\mathscr H}_x^p$ the set of all operators on $T_x$ which can be obtained from the mapping $F_x$ as above. Now note that $\mathbb H_x J$ is the matrix whose elements are the homogeneous parts of the elements of the matrix corresponding to $J$. Hence if we denote by $\handwritten{\mathscr F}_x^p$ the ideal of $\handwritten{\mathscr G}^p$ consisting of all $J\handwritten{\in}\handwritten{\mathscr G}^p$ with $F_x=0$ and $\handwritten{\xi}(x)=0$ then we obtain

\uline{Theorem 4.6.6.} $\handwritten{\mathscr H}_x^p$ is an algebra under composition of operators. The map $\mathbb H_x:\handwritten{\mathscr G}^p\handwritten{\longrightarrow}\handwritten{\mathscr H}_x^p$ is a homomorphism of algebras with kernel $\handwritten{\mathscr F}_x^p$.

Since the space $T_x$ has an inner product we have an isomorphism $T_x\handwritten{\longrightarrow}T_x^*$ which induces a natural isomorphism between $A_x^p\mathbin{\ooalign{\hfil$\mathrm{x}$\hfil\cr\hfil$\handwritten{\bigcirc}$\hfil\cr}}A_x^p$ and the space of linear transformations $A_x^p\handwritten{\longrightarrow}A_x^p$. This isomorphism provides $A_x^p\mathbin{\ooalign{\hfil$\mathrm{x}$\hfil\cr\hfil$\handwritten{\bigcirc}$\hfil\cr}}A_x^p$ with an algebraic structure. Thus $\handwritten{\mathscr S}_x^p$, the set of $C^\infty$ maps $S_x\handwritten{\longrightarrow}A_x^p\mathbin{\ooalign{\hfil$\mathrm{x}$\hfil\cr\hfil$\handwritten{\bigcirc}$\hfil\cr}}A_x^p$, is an algebra. Note that $\handwritten{\mathscr S}_x^p$ can be represented as the algebra of $\binom mp\,\mathrm{x}\,\binom mp$ matrices of $C^\infty$ functions on $S_x$.

Consider the map $\unsure{\mathbb S|S}_x:\handwritten{\mathscr H}_x^p\handwritten{\longrightarrow}\unsure{.}\handwritten{\mathscr S}_x^p$ defined by
\[
(\mathbb S_x\mathbb H_x J)(u)=\handwritten{\xi}(x)+\handwritten{\int}_{T_x} e^{iu\cdot v}\frac{F_x(v)}{\lVert v\rVert^m}\,dV_v
\tag*{(4.6.7)}
\]
\scanpage{100}
\mspage{\textcircled{26}}
Formulaire faisant intervenir $\overset{\mathbb L}{\otimes}$ et $D^-$

\textcircled{1} Généralités valables sur tout topos annelé

Module \uline{plat} $F$ : $F\otimes\cdot$ est exact.

$F$ plat $\Longleftrightarrow$ Pour tout injectif $I$, $\mathscr Hom(F,I)$ est injectif.

\uline{NB} Dans le cas actuel, il suffit de vérifier sur les fibres.

Tout \deleted{\unsure{f}}Module $F$ \deleted{admet des \unsure{ré}} est quotient d'un plat, donc admet des résolutions plates : on prend les $A_U$ ($U\in\mathrm{Ob}\,X_{\mathrm{ét}}$). Il s'ensuit \inserted{\unsure{formellt|formellmt}} que pour tt $\deleted{\unsure{K}}L_{\bullet}\in\mathrm{Ob}\,D^-(X)$, existe \unsure{liaison|homom.} $L'_{\bullet}\mathbin{-}L_{\bullet}$, avec $L'_{\bullet}$ plat en \unsure{tout degré|tous degré|tous degrés}, qui est un quasi-isomorphisme.

De ceci résulte formellement que $D^-(X)$ est équivalent : la catégorie obtenue en prenant \inserted{la catégorie $K^-(X)^{\mathrm{plat}}$} des complexes $L_{\bullet}$ bornés sup. plats \deleted{\unsure{en}} tous degrés, \deleted{\unsure{et}} en prenant les hom. à homotopie près, puis prenant un quotient par la \deleted{relation} \inserted{famille} des quasi-isomorphismes.

Ceci va nous permettre de prendre \deleted{des} ``foncteur dérivé gauche'' du foncteur $\otimes$, en prenant dans $C^-(X)^{\mathrm{plat}}$ le $\otimes$ ordinaire ; on \unsure{note|voit} que ça passe au quotient par quasi-isomorphismes. Donc si $L_{\bullet},\deleted{L}M_{\bullet}$ sont deux complexes ordinaires $\in\mathrm{Ob}\,D^-(X)$, on \deleted{définit} calcule $L_{\bullet}\overset{L}{\otimes}\deleted{L}M_{\bullet}$ en prenant des résolutions plates $L'_{\bullet}\longrightarrow{}L_{\bullet}$, $M'_{\bullet}\mathbin{-}M_{\bullet}$ et prenant $L'_{\bullet}\otimes M'_{\bullet}$.
% archive sheet number: 99
\scanpage{101}
\mspage{3.6}
% archive sheet number: 100
Corollary 3.9: Let $R\handwritten{\subset}S$ be a local algebra of the integrally closed local rings such that the quotients field extension $k\handwritten{\subset}K$ is galois. Then $r(S/R)={}$ the number of minimal primes $\underline p$ in $S^f$ such that $\underline p\cap R=0={}$ the order of the inertia group of $S$ over $R$.

\uline{Proof}: Let $U$ be the inertial ring of $S$ over $R$. Since $U$ is an unramified local $R$-algebra, we have by lemma 2.1 that $S_R^f\handwritten{\approx}S_U^f$. However $U$ is a transporter for $S$ over $R$ and hence $S_R^f\handwritten{\approx}S_U^f=S\mathbin{\ooalign{\hfil$\mathrm{x}$\hfil\cr\hfil$\handwritten{\bigcirc}$\hfil\cr}}_U S$. For a prime ideal $\underline p$ in $S_R^f$ it is clear that $\underline p\cap R=0$ if and only if $\underline p\cap U=0$, since $U$ is a localization of an integral extension of $R$. Therefore the minimal prime ideals $\underline p$ in $S_R^f$ with $\underline p\cap R=0$ are in one-one correspondence with the prime ideals in $(S\mathbin{\ooalign{\hfil$\mathrm{x}$\hfil\cr\hfil$\handwritten{\bigcirc}$\hfil\cr}}_U S)\mathbin{\ooalign{\hfil$\mathrm{x}$\hfil\cr\hfil$\handwritten{\bigcirc}$\hfil\cr}}_U K_I=K\mathbin{\ooalign{\hfil$\mathrm{x}$\hfil\cr\hfil$\handwritten{\bigcirc}$\hfil\cr}}_{K_I}K$ where $K_I$ is the field of quotients of $U$, i.e. the inertial field of $S$ over $R$. Now $K$ is a galois extension of $K_I$ with $G_I$ the inertia group of $S$ over $R$ as its galois group. We recall that there is one-one correspondence between the prime ideals in $K\mathbin{\ooalign{\hfil$\mathrm{x}$\hfil\cr\hfil$\handwritten{\bigcirc}$\hfil\cr}}_{K_I}K$ and the elements in $G_I$ under $\underline P(\handwritten{\sigma})=\mathrm{Ker}\,(\handwritten{\widetilde\sigma}=K\mathbin{\ooalign{\hfil$\mathrm{x}$\hfil\cr\hfil$\handwritten{\bigcirc}$\hfil\cr}}_{K_I}K\longrightarrow K)$ where $\handwritten{\widetilde\sigma}(x\mathbin{\ooalign{\hfil$\mathrm{x}$\hfil\cr\hfil$\handwritten{\bigcirc}$\hfil\cr}} y)=x\handwritten{\sigma}(y)$. Therefore the prime ideals $\underline p$ in $S_R^f\handwritten{\approx}S\mathbin{\ooalign{\hfil$\mathrm{x}$\hfil\cr\hfil$\handwritten{\bigcirc}$\hfil\cr}}_U S$ with $\underline p\cap R=0$ are in one-one correspondence with the elements in $G_I$ under $\underline p(\handwritten{\sigma})=\underline P(\handwritten{\sigma})\cap S_R^f=\underline P(\handwritten{\sigma})\cap(S\mathbin{\ooalign{\hfil$\mathrm{x}$\hfil\cr\hfil$\handwritten{\bigcirc}$\hfil\cr}}_U S)$. Since $G_I$ is the inertia group of $S$ over $R$, we have in particular that $\handwritten{\sigma}(S)=S$ for all $\handwritten{\sigma} \varepsilon G_I$, and hence $\underline p(\handwritten{\sigma})=\underline P(\handwritten{\sigma})\cap(S\mathbin{\ooalign{\hfil$\mathrm{x}$\hfil\cr\hfil$\handwritten{\bigcirc}$\hfil\cr}}_U S)=\mathrm{Ker}\,(\handwritten{\widetilde\sigma}:S\mathbin{\ooalign{\hfil$\mathrm{x}$\hfil\cr\hfil$\handwritten{\bigcirc}$\hfil\cr}}_U S\longrightarrow S)$ where $\handwritten{\widetilde\sigma}(x\mathbin{\ooalign{\hfil$\mathrm{x}$\hfil\cr\hfil$\handwritten{\bigcirc}$\hfil\cr}} y)=x\handwritten{\sigma}(y)$. Therefore $S_R^f/\underline p(\handwritten{\sigma})=S\mathbin{\ooalign{\hfil$\mathrm{x}$\hfil\cr\hfil$\handwritten{\bigcirc}$\hfil\cr}}_U S/\underline p(\handwritten{\sigma})=S$ and hence $e(S_R^f/\underline p(\handwritten{\sigma})\mid S)=1$ for all $\handwritten{\sigma} \varepsilon G_I$. Therefore $r(S/R)=\handwritten{\sum}_{\varepsilon G_I}e(S^f/\underline p(\handwritten{\sigma})\mid S)={}$ the number of prime ideals $\underline p$ in $S_R^f$ with $\underline p\cap R=0={}$ the order of $G_I$.

Proposition 3.10: Let $R$ be an integrally closed local domain with quotient field $k$. Let $K$ be a finite algebraic extension of $k$ and $S$ the
\scanpage{102}
Les \deleted{objets} faisceaux de cohomologie de $L_{\bullet}\overset{L}{\otimes}M_{\bullet}$ sont appelés \uline{hypertor} de $L_{\bullet},M_{\bullet}$, notation $\mathbb T\mathrm{or}_i^{\underline O}(L_{\bullet},M_{\bullet})$. \deleted{\unsure{Si $L_{\bullet}$ ou $M$|Si $L_{\bullet},M$}} Comme il s'agit ici de foncteurs \inserted{``exacts''} de catégories triangulées, tout triangle exact en l'un des arguments donne lieu à la suite exacte infinie habituelle. En particulier, les $\mathbb T\mathrm{or}_i(\ ,M_{\bullet})$ définissent un foncteur cohomologique en le premier argument, lorsqu'on prend pour celui-ci un Module ordinaire.

La formation du $\deleted{\unsure{K}}L_{\bullet}\overset{\mathbb L}{\otimes}M_{\bullet}$ commute au changement de base \uline{plat} \inserted{$f:X\longrightarrow Y$}, quand celui-ci est défini \unsure{naïvement} : \deleted{\unsure{[mais alors}}
\[
\boxed{\mathbb L f^*(M_{\bullet}\overset{\mathbb L}{\otimes}L_{\bullet})\simeq\mathbb L f^*(M_{\bullet})\overset{\mathbb L}{\otimes}\mathbb L f^*(L_{\bullet})}
\]
Pour définir le foncteur \unsure{$f\mathbin{-}$|$f^*$ ou} $\mathbb L f^*$ entre $D^-$, il faut noter que la \uline{platitude} de $f$ implique que $f^*$ transforme quasi-isomorphismes en quasi-isomorphismes. D'autre part, $f^*$ transforme plats en plats, d'où la commutation.
% slanted
\margin{est-ce vrai dans le contexte des topos généraux ?}
--- En particulier, la formation des $L_{\bullet}\overset{\mathbb L}{\otimes}M_{\bullet}$ commute au passage aux fibres, ce qui dans le cas qui nous occupe permettra \inserted{\unsure{pfois|parfois|tjours}} un calcul pratiquement trivial des $\overset{\mathbb L}{\otimes}$.
% archive sheet number: 101
\scanpage{103}
\mspage{3.5}
% archive sheet number: 102
follows from lemma 3.4 that $e(\overset{\handwritten{\sim}}{S}/\underline p\mid D)$ is an integer for all $\underline p$ with $\underline p\cap R=0$ and consequently $e(S/R)=e(\overset{\handwritten{\sim}}{S}/D)=\handwritten{\sum}_{\underline p}L(\overset{\handwritten{\sim}}{S}\underline p)e(\overset{\handwritten{\sim}}{S}/\underline p\mid D)$ is an integer.

Theorem 3.8: Let $R\handwritten{\subset}S$ be a local algebra of local domains with quotient fields $k\handwritten{\subset}K$ and assume that $S$ is integrally closed. Then $r(S/R)=\handwritten{\sum}_{\underline p}L(S^f_{\underline p})e(S^f/\underline p\mid S)$ and $e(S^f/\underline p\mid S)$ are all nonzero integers, where $\underline p$ ranges through the minimal primes in $S^f$ such that $\underline p\cap R=0$. In the above formula, $L(S^f_{\underline p})=1$ for all $\underline p$ with $\underline p\cap R=0$ if and only if $K$ is separably algebraic over $k$.

\uline{Proof}: From lemma 3.6, we have $r(S/R)=e(S^f/S)=\handwritten{\sum}_{\underline p}L(S^f_{\underline p})e(S^f/\underline p\mid S)$ where $\underline p$ ranges through the minimal prime ideals in $S^f$ such that $\dim S^f/\underline p=\dim S^f$. Put $\underline p_0=\mathrm{Ker}(S^f\longrightarrow S)$. Then $\underline p_0$ is a prime ideal in $S^f$ such that $\underline p_0\cap S=0$. Therefore we are in the position to apply lemma 3.5 since $S^f$ is a local $S$-algebra and $S$ is integrally closed. Therefore $\dim S^f/\underline p=\dim S^f$ if and only if $\underline p\cap S=0$. However, it is clear that $\underline p\cap S=0$ if and only if $\underline p\cap R=0$ since $S$ is a localization of an integral extension of $R$. Since $S$ is integrally closed and $S^f/\underline p$ and $S$ have the same residue class field, it follows from proposition 3.7 that $e(S^f/\underline p\mid S\handwritten{)}\unsure{'}$ is an integer. This proves the first assertion. Now it is clear that $S^f_{\underline p_0}=K^f_k$ and hence $L(S^f_{\underline p_0})=L(K^f_k)=[K:k]_i$ by corollary 2.9. Thus if $L(S^f_{\underline p})=1$ for all $\underline p$ with $\underline p\cap R=0$, then in particular we must have that $L(S^f_{\underline p_0})=[K:k]_i=1$, i.e. $K$ is separably algebraic over $k$. Conversely if $K$ is separably algebraic over $k$ then $K\mathbin{\ooalign{\hfil$\mathrm{x}$\hfil\cr\hfil$\handwritten{\bigcirc}$\hfil\cr}}_k K$ is a separable $k$-algebra and consequently $(K\mathbin{\ooalign{\hfil$\mathrm{x}$\hfil\cr\hfil$\handwritten{\bigcirc}$\hfil\cr}}_k K)_{\underline P}$ is a field for all prime ideal $\underline P$ in $K\mathbin{\ooalign{\hfil$\mathrm{x}$\hfil\cr\hfil$\handwritten{\bigcirc}$\hfil\cr}}_k K$. Consequently $L(S^f_{\underline p})=1$ for all $\underline p$ with $\underline p\cap R=0$.
\scanpage{104}
Variante globale : on définit des hypertor globaux \deleted{\unsure{par}} prenant $\mathbb R\Gamma_X\,L_{\bullet}\overset{\mathbb L}{\otimes}M_{\bullet}$.
% slanted
\margin{Mais il faut supposer, si on ne fait pas d'\unsure{hyp.} suppl. \unsure{pour|sur} $L_{\bullet},M_{\bullet}$, que $\Gamma_X$ est de dim \unsure{coh} finie !}

\uline{Lois formelles les plus évidentes} (lois internes).

1) Associativité de $\overset{\mathbb L}{\otimes}$

\unsure{2}) Commutativité (attention au signe !)

3)
% slanted
\margin{N.B. $\mathbb R\mathscr Hom(L_{\bullet},K^{\bullet})$ défini dès que $K^{\bullet}\in\mathrm{Ob}\,D^+(X)$, mais $L_{\bullet}$ peut être quelconque}
\diagram{\begin{tikzcd}[ampersand replacement=\&,column sep=.2em,row sep=.4em]
\mathbb R\mathscr Hom(L_{\bullet}\overset{L}{\otimes}M_{\bullet},K^{\bullet}) \& {}\simeq\mathbb R\mathscr Hom(L_{\bullet},\mathbb R\mathscr Hom(M_{\bullet},K^{\bullet}))\\
\& {}\simeq\mathbb R\mathscr Hom(M_{\bullet},\mathbb R\mathscr Hom(L_{\bullet},K^{\bullet}))\\
\mathbb R\mathrm{Hom}(L_{\bullet}\overset{\mathbb L}{\otimes}M_{\bullet},K^{\bullet})\arrow[uu,out=165,in=165] \& {}\simeq\mathbb R\mathrm{Hom}(L_{\bullet},\mathbb R\deleted{\unsure{H}}\mathscr Hom(M_{\bullet},K^{\bullet}))\\
\& {}\simeq\mathbb R\mathrm{Hom}(M_{\bullet},\mathbb R\mathscr Hom(L_{\bullet},K^{\bullet}))
\end{tikzcd}}
\[
\begin{aligned}
&\deleted{\mathrm{Hom}(L_{\bullet},\mathbb R\mathscr Hom}\\
\mathrm{Hom}(L_{\bullet}\overset{\mathbb L}{\otimes}M_{\bullet},K^{\bullet})&\simeq\mathrm{Hom}(L_{\bullet},\mathbb R\mathscr Hom(M_{\bullet},K^{\bullet})\\
&\simeq\mathrm{Hom}(M_{\bullet},\mathbb R\mathscr Hom(L_{\bullet},K^{\bullet}))
\end{aligned}
\]

\uline{Remarque sur changement de base non plat}

On définit $\mathbb L f^*:\deleted{\unsure{D}}D^-(Y)\longrightarrow D^-(X)$ encore en passant par $K^-(Y)^{\mathrm{plat}}/\text{quasi-isom}$, en notant que $f^*$ \deleted{\unsure{sur}} \deleted{\unsure{$K^-$}} envoie $K^-(Y)^{\unsure{\mathrm{plat}}}$ dans $K^-(X)^{\mathrm{plat}}$ et transforme qu.\ isom en idem. Avec cette définition, on aura encore
\[
\mathbb L f^*(L_{\bullet}\overset{\mathbb L}{\otimes}M_{\bullet})\simeq\mathbb L f^*(L_{\bullet})\overset{\mathbb L}{\otimes}\mathbb L f^*(M_{\bullet})
\]
Notons aussi formule de transitivité
% archive sheet number: 103
\scanpage{105}
\mspage{3.4}
% archive sheet number: 104
unramified over $S$. Since $\overline D=\overline R$, we have a natural epimorphism
$D\mathbin{\ooalign{\hfil$\mathrm{x}$\hfil\cr\hfil$\handwritten{\bigcirc}$\hfil\cr}}_R S\xrightarrow{\handwritten{\varphi}}\overline S$. Let $\underline M=\mathrm{Ker}\,\handwritten{\varphi}$ and set $\overset{\handwritten{\sim}}{S}=(D\mathbin{\ooalign{\hfil$\mathrm{x}$\hfil\cr\hfil$\handwritten{\bigcirc}$\hfil\cr}}_R S)_{\underline M}$. Then $\overset{\handwritten{\sim}}{S}$ is flat and
unramified over $S$. Therefore $e(S/R)=e(\overset{\handwritten{\sim}}{S}/D)$ by lemma 2.3. From
lemma 3.6, we get $e(\overset{\handwritten{\sim}}{S}/D)=\handwritten{\sum}_{\underline p}L(\overset{\handwritten{\sim}}{S}_{\underline p})e(\overset{\handwritten{\sim}}{S}/\underline p\mid D)$ where $\underline p$ ranges through the
minimal primes in $\overset{\handwritten{\sim}}{S}$ such that $\dim \overset{\handwritten{\sim}}{S}/\underline p=\dim \overset{\handwritten{\sim}}{S}$. Now if we let $\underline p_0={}$
$\mathrm{Ker}(D\mathbin{\ooalign{\hfil$\mathrm{x}$\hfil\cr\hfil$\handwritten{\bigcirc}$\hfil\cr}}_R S\longrightarrow DS)$, then $\underline p_0\cap R=0$ and thus we are in a position to
apply lemma 4.5. Thus we have that $e(\overset{\handwritten{\sim}}{S}/D)=\handwritten{\sum}_{\underline p}L(\overset{\handwritten{\sim}}{S}_{\underline p})e(\overset{\handwritten{\sim}}{S}/\underline p\mid D)$ where
$\underline p\cap R=0$. We claim that the local $D$-algebras $\overset{\handwritten{\sim}}{S}/\underline p$ are all finite $D$-module.
For, since there is one-one correspondence between the prime ideals $\underline p$ in
$D\mathbin{\ooalign{\hfil$\mathrm{x}$\hfil\cr\hfil$\handwritten{\bigcirc}$\hfil\cr}}_R S$ with $\underline p\cap R=0$ and the prime ideals in $(D\mathbin{\ooalign{\hfil$\mathrm{x}$\hfil\cr\hfil$\handwritten{\bigcirc}$\hfil\cr}}_R S)\mathbin{\ooalign{\hfil$\mathrm{x}$\hfil\cr\hfil$\handwritten{\bigcirc}$\hfil\cr}}_R k=L'_D\mathbin{\ooalign{\hfil$\mathrm{x}$\hfil\cr\hfil$\handwritten{\bigcirc}$\hfil\cr}}_k K$
(where $L'_D$ is the quotient field of $D$), we see that $\overset{\handwritten{\sim}}{S}/\underline p$ is isomorphic to (and
hence we shall identify) a localization of $\handwritten{\sigma}(S)D$ as $D$-algebras, for some
conjugation $\handwritten{\sigma}$ of $S$ over\unsure{'} $R$. (The reader may observe that $\handwritten{\sigma}(S)D$ corresponding
to $\overset{\handwritten{\sim}}{S}/\underline p$ has, by the definition of $\overset{\handwritten{\sim}}{S}=(D\mathbin{\ooalign{\hfil$\mathrm{x}$\hfil\cr\hfil$\handwritten{\bigcirc}$\hfil\cr}}_R S)_{\underline M}$, a maximal ideal which
dominates both the maximal ideals of $S^{\handwritten{\sigma}}$ and $D$). The quotient field of
$\overset{\handwritten{\sim}}{S}/\underline p$ is $L'_D\handwritten{\sigma}(K)$ and $L'_D\handwritten{\sigma}(K')$ is the separable closure of $L'_D$ in
$L'_D\handwritten{\sigma}(K)$. Since $\overset{\handwritten{\sim}}{S}/\underline p$ is a local $D$-algebra, we can find a local $D$-algebra
$A$ with quotient field $L'_D\handwritten{\sigma}(K')$ such that $D\handwritten{\subset} A\handwritten{\subset} \overset{\handwritten{\sim}}{S}/\underline p$. Since $L'_D\handwritten{\sigma}(K)$
is purely inseparable over $L'_D\handwritten{\sigma}(K')$, every integral extension of a local
ring $A$ contained in $L'_D\handwritten{\sigma}(K)$ is local and consequently local $A$-algebra
$\overset{\handwritten{\sim}}{S}/\underline p$ is a finite $A$-module. On the other hand, since the integral closure
of $D$ in $L'$ is a finite local $D$-algebra (proposition 1.1), every local $D$-algebra
contained in $L'$ is a finite local $D$-algebra. In particular, the
local $D$-algebra $A$, being contained in $L'_D\handwritten{\sigma}(K')\handwritten{\subset} L$, is a finite local
$D$-algebra. Therefore we conclude that $\overset{\handwritten{\sim}}{S}/\underline p$ is a finite local $D$-algebra
for all $\underline p$ with $\underline p\cap R=0$. Furthermore the residue field of $S/\underline p$, being
isomorphic to $\overline S$, is a simple extension of $\overline D=\overline R$ by the hypothesis. It
\scanpage{106}
% archive sheet number: 105
\[
\mathbb L(gf)^* \simeq \mathbb L f^*\,\mathbb L g^*
\]
valable sans hypothèse sur $f,g$ composables.

Écriture \deleted{\unsure{$L$|$\otimes$}} (Si
\diagram{\begin{tikzcd}[ampersand replacement=\&,column sep=small]
f:\underset{L_\bullet}{X}\arrow[r]\&\underset{M_\bullet}{Y}\&L_\bullet\overset{\mathbb L}{\otimes}_Y M_\bullet \simeq L_\bullet\overset{\mathbb L}{\otimes}\mathbb L f^*(M_\bullet)
\end{tikzcd}}

\uline{Formule d'adjonction} pour $f:X\longrightarrow Y$ :

\diagram{\begin{tikzcd}[ampersand replacement=\&,column sep=small]
X\arrow[d,"f"']\&K^\bullet\\
Y\&L_\bullet
\end{tikzcd}}
\deleted{\unsure{$K$}}
\[
\boxed{\operatorname{Hom}(\mathbb L f^*(L_\bullet),K^\bullet)\xrightarrow{\sim}
\operatorname{Hom}(\deleted{\unsure{f_*}}L_\bullet\unsure{)},\mathbb R f_*(K^\bullet))}
\]

En effet, on peut supposer $K^\bullet$ \deleted{\unsure{p}} injectif, $L_\bullet$ plat, alors le premier mb est \deleted{(\unsure{comme} $K^\bullet$ injectif) $H^0\,\operatorname{Hom}^\bullet(\unsure{R} f^*(L_\bullet),K^\bullet)$} le $H^0$ de
\[
\operatorname{Hom}^\bullet(f^*L_\bullet,K^\bullet)\simeq
\operatorname{Hom}^\bullet(L_\bullet,f_*(K^\bullet))\qquad \text{OK.}
\]

\unsure{Variante de projection|Variante ou projection|Variante ou généralisation}
\[
\mathbb R f_*\mathbb R\mathscr H\!om^\bullet(\mathbb L f^*L_\bullet,K^\bullet)
\simeq\mathbb R\mathscr H\!om^\bullet(L_\bullet,\mathbb R f_*(K^\bullet))
\]

\[
\begin{array}{|l}
\text{\uline{Cor}}\qquad \mathbb R\operatorname{Hom}^\bullet(\mathbb L f^*(L_\bullet),K^\bullet)\simeq
\mathbb R\operatorname{Hom}^\bullet(L_\bullet,\mathbb R f_*(K^\bullet))\\\hline
\end{array}
\]

% slanted
\margin{Attention, on utilise que si $f$ est plat alors $f_*(K^\bullet)$ injectif si $K^\bullet$ injectif. Dans le cas général, par dévissage sur $L_\bullet$ se ramener au cas où $L=\deleted{\unsure{L}}A_{Y',Y}^{\unsure{!|g}}$ et le vérifier directement dans ce cas :}
\[
\begin{aligned}
&\mathbb R\operatorname{Hom}^\bullet(f^*L_\bullet,K^\bullet)\\
&=\mathbb R\operatorname{Hom}^\bullet(A_{X'},K^\bullet)
=\deleted{\unsure{\operatorname{Hom}}}\mathbb R\Gamma_{X'}\,\deleted{\unsure{i^*}}\underbrace{\mathbb L j^*(K^\bullet)}_{\text{image sur }X'}
\simeq \mathbb R\Gamma_{Y'}(\underbrace{\mathbb R f'_*\mathbb L j^*(K^\bullet)}_{\mathbb L i^*(\mathbb R f_*(K^\bullet))})\\
&\simeq\mathbb R\operatorname{Hom}^\bullet(A_{Y'},\mathbb L i^*(\mathbb R f_*(K^\bullet)))\\
&=\mathbb R\operatorname{Hom}^\bullet(L_\bullet,\mathbb R f_*(K^\bullet))
\end{aligned}
\]
\margin{car $j:Y'\to Y$ est un morphisme de localisation \ldots}
\diagram{\begin{tikzcd}[ampersand replacement=\&]
X'\arrow[r,"j"]\arrow[d,"f'"']\&X\arrow[d,"f"]\\
Y'\arrow[r,"i"']\&Y
\end{tikzcd}}
\scanpage{107}
\mspage{3.3}
% archive sheet number: 106
Lemma 3.5: Let $S$ be a local algebra over an integrally closed local
domain $R$. If $S$ contains a prime ideal $\underline p$ such that $\underline p\cap R=0$, then
$\dim S=\dim R$ and, for each prime ideal $\underline p$ in $S$, we have that $\dim S/\underline p=$
$\dim S$ if and only if $\underline p\cap R=0$.

\uline{Proof}: Let $\underline p$ be a prime ideal in $S$ such that $\underline p\cap R=0$. Then
$\dim S/\underline p=\dim R$ by "going down theorem" since $R\handwritten{\subset} S/\underline p$ is a local algebra
and $R$ is integrally closed. In particular, we have $\dim S\geq\dim R$ and
consequently $\dim S=\dim R$ since $S$ is a local $R$-algebra. For a prime ideal
$\underline p$ in $S$, $\dim S/\underline p\leq\dim R/\underline p\cap R$ since $S/\underline p$ is a local algebra over $R/\underline p\cap R$.
Therefore $\dim S/\underline p=\dim S(=\dim R)$ entails that $\underline p\cap R=0$ since $R$ is a
domain. Conversely, if $\underline p\cap R=0$, we have already observed that $\dim S/\underline p=$
$\dim R$, i.e. $\dim S/\underline p=\dim S$.\hfill\unsure{;|:}

The following lemma is well-known ([11]):

Lemma 3.6: Let $R$ be a local ring and $E$ a finitely generated $R$-module.
Then for each ideal of definition $\underline q$ in $R$, we have $e_E(\underline q)\unsure{.|,}=
\handwritten{\sum}_{\underline p} L_{R_{\underline p}}(E_{\underline p})e_{R/\underline p}(\underline q)$
where $\underline p$ runs through all minimal primes of $R$ such that $\dim R/\underline p=\dim R$.

Proposition 3.7: Let $R\handwritten{\subset} S$ be a local algebra of local domains. If
$R$ is integrally closed and $\bar S$ is a simple extension of $\bar R$, then $e(S/R)$ is
an integer.

\uline{Proof}: Let $k\handwritten{\subset} K$ be the quotient field extension of $R\handwritten{\subset} S$ and let $L$
be the normal closure of $K$ over $k$. Let $L'$, $K'$ be the separable closure
of $k$ in $L$, $K$ respectively. Then $L'$ contains $K'$ and is galois over $k$.
Since $R\handwritten{\subset} S$ is a local algebra, we can find $S'$ such that $R\handwritten{\subset} S'\handwritten{\subset} S$ and $S'$
is a local $R$-algebra with quotient field $K'$. Let $T$ be a localization by a
maximal ideal of the integral closure of $R$ in $L'$ such that $T$ dominates $S'$,
and let $D$ be the decomposition ring of $T$ over $R$. Then $D$ is flat and
unramified over $R$ by proposition 1.1 and hence $D\mathbin{\ooalign{\hfil$\mathrm{x}$\hfil\cr\hfil$\handwritten{\bigcirc}$\hfil\cr}}_R S$ is flat and
\scanpage{108}
% archive sheet number: 107
\uline{Remarque.} Nous nous sommes bornés ici aux formules générales valables sans conditions de régularité (genre pseudo-cohérence, ou quasi-cohérence) qui donnent de nombreuses autres formules. D'autre part, nous nous sommes bornés aux cas d'existence \inserted{“naturelle” i.e. générale} des op. envisagées $\overset{\mathbb L}{\otimes}$, \inserted{$\mathbb R\mathscr H\!om$}, $\mathbb L f^*$, $\mathbb R f_*$, sans examiner les variantes \unsure{pr.|pro}\deleted{\unsure{venant|venant de}} d'hyp. de Tor-dim ou Ext-dim. finie \ldots.

\begin{center}\rule{0.95cm}{0.4pt}\\[-0.75em]\rule{0.72cm}{0.4pt}\end{center}

\textcircled{2} Formules \unsure{sur les précédentes|sur les précédents|pour les précédentes|pour les précédents}, faisant intervenir égalemt $!$ en plus de $\overset{\mathbb L}{\otimes}$ : Künneth.

\deleted{\uline{\unsure{Théorème}}. Soit $f:X\longrightarrow Y$ un morphisme compactifiable.}
\deleted{Nous \unsure{pouvons|pourr}}

Notons que en général nous sommes obligés de travailler avec des faisceaux d'anneaux constant $\mathbb Z/n\mathbb Z=A$ qui n'est pas un corps ni de tor-dimension \inserted{globale} finie, donc dès qu'on commence à faire du $\overset{\mathbb L}{\otimes}$, \deleted{\unsure{en}} en \unsure{partant|faisant} avec un faisceau réduit au degré 0, on va \unsure{tomber|déborder} \unsure{dans|des} $D^-$. Il faut
\scanpage{109}
\mspage{3.2}
% archive sheet number: 108
\uline{Proof}: 1) Assume that $S$ is finitely generated as $R\unsure{'}$-module. Then
$\mathrm{w.}\ \mathrm{hd}_R S<\infty$ implies $\mathrm{hd}_R S<\infty$. Let $0\longrightarrow X_t\longrightarrow\ldots\longrightarrow X_1\longrightarrow X_0\longrightarrow
S\longrightarrow 0$ be an acyclic resolution of $S$ by $R$-free modules of finite rank,
and let $\underline m$ be the maximal ideal of $R$. Then by the additivity of the multiplicity,
we have $\dfrac{e_S(\underline m)}{e_R(\underline m)}=\handwritten{\sum}_i(-1)^i\,e_{X_i}(\underline m)/e_R(\underline m)=\handwritten{\sum}_i(-1)^i[X_i:R]=$
\[
\begin{aligned}
&\handwritten{\sum}_i(-1)^i[X_i\mathbin{\ooalign{\hfil$\mathrm{x}$\hfil\cr\hfil$\handwritten{\bigcirc}$\hfil\cr}}_R\bar R:R]=\handwritten{\sum}_i(-1)^i[\operatorname{Tor}_i^R(S,\bar R):\bar R]=\qquad\unsure{:}\\
&[\bar S:\bar R]\cdot[\handwritten{\sum}_i(-1)^i L_S(\operatorname{Tor}_i^R(S,\bar R))].
\end{aligned}
\]
Therefore $e(S/R)=\dfrac{e_S(\underline m S)}{e_R(\underline m)}=$
\[
\dfrac{1}{[\bar S:\bar R]}\cdot\dfrac{e_S(\underline m)}{e_R(\underline m)}=\handwritten{\sum}_i(-1)^i L_S(\operatorname{Tor}_i^R(S,\bar R)).
\]

2) General case: Since $\bar S$ is finite dimensional over $\bar R$, $\overset{\handwritten{\scriptscriptstyle\wedge}}{S}$ is finitely
generated as $\overset{\handwritten{\scriptscriptstyle\wedge}}{R}$-module. Furthermore, by lemma 3.1, we have $\mathrm{hd}_{\overset{\handwritten{\scriptscriptstyle\wedge}}{R}}\overset{\handwritten{\scriptscriptstyle\wedge}}{S}<\infty$.
Therefore we are brought to the above situation and hence $e(S/R)=e(\overset{\handwritten{\scriptscriptstyle\wedge}}{S}/\overset{\handwritten{\scriptscriptstyle\wedge}}{R})=$
\[
\handwritten{\sum}_i(-1)^i L_{\overset{\handwritten{\scriptscriptstyle\wedge}}{S}}(\operatorname{Tor}_i^{\overset{\handwritten{\scriptscriptstyle\wedge}}{R}}(\overset{\handwritten{\scriptscriptstyle\wedge}}{S},\bar R))=\handwritten{\sum}_i(-1)^i L_S(\operatorname{Tor}_i^R(S,\bar R)).
\]

Corollary 3.3: Let $S$ be a local $R$-algebra over a local ring $R$. If
$\mathrm{w.}\ \mathrm{hd}_S S^f<\infty$, then $r(S/R)$ is an integer, equal to $\handwritten{\sum}_i(-1)^i L_{S^f}(\operatorname{Tor}_i^S(S^f\underline m\,\bar S))=$
\[
\handwritten{\sum}_i(-1)^i L_{S^f}(S^f\mathbin{\ooalign{\hfil$\mathrm{x}$\hfil\cr\hfil$\handwritten{\bigcirc}$\hfil\cr}}_{S^e}\operatorname{Tor}_i^S(S^e,\bar S)).
\]

Lemma 3.4: Let $R\handwritten{\subset} S$ be local domains such that $S$ is finitely
generated as $R$-module. If $R$ is integrally closed and $\bar S$ is a simple
extension of $\bar R$, then $e(S/R)$ is an integer.

\uline{Proof}: Let $\alpha$ be $\underset{\unsure{:}}{\text{an}}$ element in $S$ such that $\bar S=\bar R(\bar\alpha)$ where $\bar\alpha$ is the
image of $\alpha$ under the canonical map $S\longrightarrow\bar S$. The subring $R[\alpha]$ is $R$-free
since $R$ is integrally closed, and is a local $R$-algebra since $S$ is a finite
local $R$-algebra. Therefore $e(R[\alpha]/R)=[R[\alpha]:R]/[\bar R[\bar\alpha]:\bar R]=$
$[R[\alpha]:R]/[\bar S:\bar R]$ is an integer by proposition 3.2. Consequently
$e(S/R)=[S:R]/[\bar S:\bar R]=[S:R[\alpha]]\cdot e(R[\alpha]/R)$ is an integer.
\scanpage{110}
% archive sheet number: 109
donc, pour cette raison, savoir définir les $\mathbb R f_*$, $\mathbb R_! f_*$ par des arguments qui ne sont pas \unsure{des|du|dans} $D^+$, \inserted{comme foncteurs $D^-(X)$ --- $D^-(Y)$, $D(X)$ --- $D(Y)$.} Pour $\mathbb R f_*$, ce sera possible dès que $f_*$ est de dim cohomologique finie. C'est peut-être vrai dès que $Y$ est noethérien, cela \deleted{\unsure{est}} \inserted{résulte} en tous cas de la \unsure{finitude} de la dim ch. si $Y$ est, \deleted{soit de} soit de type fini sur un corps, soit de type fini sur $\operatorname{Spec}\mathbb Z$. Pour $\mathbb R_! f_*$, on est en meilleure posture encore, grâce au changt de base propre, donc \uline{si} $Y$ \uline{qu. cpct}. on a
\diagram{\[
\left.
\begin{aligned}
\mathbb R_! f_*:\quad&D(X)\;\text{---}\;D(Y)\\
&D^+(X)\;\text{---}\;D^+(Y)\\
&D^-(X)\;\text{---}\;D^-(Y)\\
&D^b(X)\;\text{---}\;D^b(Y)
\end{aligned}
\right\}\quad\begin{gathered}\text{si }Y\text{ qu. cpct}\\ f\text{ compactifiable}\end{gathered}
\]}

Comme on y est, notons
\diagram{\[
\left.
\begin{aligned}
\mathbb R f_*:\quad&D(X)\;\text{---}\;D(Y)\\
&D^{\unsure{+|e}}(X)\;\text{---}\;D^{\unsure{+|c}}(Y)\\
&D^-(X)\;\text{---}\;D^-(Y)\\
&D^b(X)\;\text{---}\;D^b(Y)
\end{aligned}
\right\}\quad\begin{gathered}\text{si }Y\text{ de t.f.}\\\text{sur un corps}\\\text{ou sur }\operatorname{Spec}\mathbb Z,\\\text{et }f\text{ de t.f.}\end{gathered}
\]}
\diagram{\[
\begin{aligned}
\mathbb L f^*:\quad&D(\deleted{\unsure{X}}Y)\;\text{---}\;D(X)\\
&D^-(Y)\;\text{---}\;D^-(X)\\
&D^+(Y)\;\text{---}\;D^+(X)\qquad D^b(Y)\to D^b(X)
\end{aligned}
\]}
